% Kapitel 7: Qualitätssicherung und Tests
\section{Qualitätssicherung und Tests}
\label{sec:qualitaetssicherung}

Das folgende Kapitel beschreibt die Maßnahmen zur Sicherstellung der Softwarequalität sowie die Bewertung der Qualität des Gesamtsystems unter realen Einsatzbedingungen.


\subsection{Teststrategie}

Die Teststrategie des Projekts kombiniert Unit-Tests zur Überprüfung einzelner Komponenten mit Integrationstests, welche das Zusammenspiel der Services auf API-Ebene validieren.
Als Test-Framework wird Pytest \cite{pytest} eingesetzt, ergänzt durch das Pytest-Coverage-Modul \cite{pytest_cov} für Codeabdeckungsanalysen.
Hierbei wurde das Ziel festgelegt, eine Testabdeckung von über 90\% für jeden Python-Service zu erreichen. \\
Dies konnte in der Implementierung mit einer Testabdeckung von 91\% für den Data-Collection Service und 94\% für den Notification Service auch erreicht werden.


\subsubsection{Unit-Tests}

Die Unit-Tests decken die Handler-Klassen beider Python-Services ab und sind nach dem Arrange-Act-Assert-Pattern (AAA) strukturiert.
Externe Abhängigkeiten wie die Synology-API, Plate Recognizer und der SMTP-Server werden durch die Mocking-Library \texttt{unittest.mock} simuliert, um die Tests unabhängig von externen Services ausführbar zu machen.

Im Data Collection Service werden alle vier Handler-Klassen einzeln getestet:

\begin{itemize}
    \item \textbf{CameraHandler:}
          Tests, welche die Interaktion mit der Synology-API durch gemockte HTTP-Responses simulieren.
          Dies umfasst Erfolgs- und Fehlerszenarien für Authentifizierung, Kameraabruf und Snapshot-Erstellung.

    \item \textbf{PlateRecognizerHandler:}
          Tests für den ALPR-API-Aufruf, geprüft mit gemockten HTTP-Responses für erfolgreiche und fehlgeschlagene Anfragen.

    \item \textbf{CountryHandler:}
          Die Bezirkserkennung für österreichische und slowenische Kennzeichen mit verschiedenen Eingabekombinationen wird mit Tests überprüft (z.B. unbekannter Ländercode, ein- und zweibuchstabige Kürzel).

    \item \textbf{DatabaseHandler:}
          Tests für Datenextraktion aus API-Antworten, Hash-Normalisierung (Groß-/Kleinschreibung, Leerzeichen) und die zeitfensterbasierte Duplikatserkennung mit Datenbankinteraktionen.
\end{itemize}

\newpage
Im Notification Service werden ebenfalls die Handler und Datenbankschicht getestet:

\begin{itemize}
    \item \textbf{EmailHandler:}
          Tests für den SMTP-Versand mit gemocktem SMTP-Server, Bulk-Versand mit teilweisen Fehlern sowie der korrekte Versand von Alert- und Update-E-Mails.

    \item \textbf{DatabaseHandler:}
          CRUD-Operationen für Benutzerpräferenzen, einschließlich Duplikat-E-Mail-Erkennung und Filterung nach Benachrichtigungspräferenzen.
\end{itemize}


\subsubsection{Integrationstests}

Die Integrationstests überprüfen die vollständigen API-Endpunkte unter Verwendung von FastAPI.
Dieses ermöglicht das Senden von HTTP-Requests gegen die Applikation, ohne einen externen Server starten zu müssen.

Für den Notification Service umfasst dies die vollständige CRUD-Schnittstelle der Benutzerpräferenzen (Create, Read, Update, Delete) sowie die Benachrichtigungs-API mit verschiedenen Empfänger-Szenarien (alle Abonnenten, spezifische Empfänger, keine Empfänger, Zustellfehler).


\subsubsection{Testisolation und Datenbankstrategie}

Die Tests werden innerhalb der CI-Pipeline mit den zugehörigen Service-Containern gegen den PostgreSQL-Service ausgeführt, anstatt eine In-Memory-Datenbank wie SQLite als Ersatz zu verwenden.
Diese Entscheidung wurde getroffen, da SQLite wesentliche benutzte PostgreSQL-Funktionalitäten nicht identisch abbilden kann.
Beispiele hierfür sind die Verwendung des nativen PostgreSQL-Enum-Datentyps \cite{postgresql_enum} für die Fahrzeugorientierung sowie die Verwendung von zeitzonenbehafteten Zeitstempeln \cite{postgresql_datetime}.
SQLite bietet für diese Typen keine native Unterstützung \cite{sqlite_datatypes}.
Durch die Nutzung derselben Datenbank-Engine in den Tests wird sichergestellt, dass die Tests tatsächlich das Verhalten der Produktionsumgebung widerspiegeln und keine falsch-positiven Ergebnisse durch abweichendes Datenbankverhalten entstehen.

\newpage
Beide dieser Services verwenden ein identisches Muster zur Datenisolation in ihrer Pytest-Konfigurationsdatei:

\begin{lstlisting}[language=Python, caption={Pytest-Fixture für Datenisolation}, label={lst:conftest}]
@pytest.fixture(scope='function')
def db():
    """
    Provides a SQLAlchemy session for each test function.
    Each test will run within its own transaction, which is 
    then rolled back at the end of the test, ensuring data 
    isolation without dropping tables.
    """
    connection = test_engine.connect()
    transaction = connection.begin()
    session = TestingSessionLocal(bind=connection)

    try:
        yield session
    finally:
        session.close()
        transaction.rollback()
        connection.close()
\end{lstlisting}

Jeder Test wird innerhalb einer eigenen Datenbank-Transaktion ausgeführt, die nach dem Test automatisch zurückgerollt wird.
Dadurch werden alle Datenbankänderungen eines Tests nach dessen Abschluss verworfen und die Tests beeinflussen sich gegenseitig nicht.
Die Reihenfolge der Testausführung ist somit beliebig.

Für die Integrationstests wird zusätzlich die FastAPI Dependency-Injection überschrieben, sodass der Test-Client die Test-Datenbanksession und, im Fall des Notification Service, einen Test-API-Key verwendet.

\newpage
\subsection{Qualitätssicherung der Kennzeichenerkennung}
\label{sec:qs_alpr}

Die Bewertung der Erkennungsrate in einem ALPR-System wird dadurch erschwert, dass eine exakte Dunkelziffer, also der Anteil der Fahrzeuge, welche das System passieren aber nicht erkannt werden, ohne eine unabhängige Referenzmessung nicht ermittelt werden kann.
Eine systematische Gegenüberstellung jeder einzelnen physischen Durchfahrt mit der entsprechenden Systemerkennung wurde im Rahmen dieser Arbeit nicht durchgeführt.
Dies begründet sich dadurch, dass eine bloße Erfassung der Fahrzeuganzahl ohne erfolgreiche Kennzeichenidentifikation einen geringen operativen Zusatznutzen bieten würde, da die eigentliche Relevanz der Daten in der Analyse von Trends und statistischen Mittelwerten liegt.
Die folgenden Erkenntnisse basieren daher auf Beobachtungen während des Betriebszeitraums sowie auf der Analyse der gespeicherten Erkennungsdaten.

\subsubsection{Herausforderungen und Umgebungseinflüsse}

Mehrere Umgebungsfaktoren beeinflussen die Erkennungsqualität im täglichen Betrieb maßgeblich:

\begin{itemize}
    \item \textbf{Lichtverhältnisse:} Da die Zufahrt nach Süden ausgerichtet ist, ergeben sich bei tiefem Sonnenstand (insbesondere in den Morgen- und Abendstunden sowie im Winter) Gegenlichtbedingungen, welche die Kennzeichenlesbarkeit beeinträchtigen können.
    \item \textbf{Witterungseinflüsse:} Bei starkem Regen, Schneefall oder Nebel ist mit einer reduzierten Erkennungsrate zu rechnen, da die Sichtbarkeit der Kennzeichen physisch eingeschränkt ist.
    \item \textbf{Nachtbetrieb:} Nächtliche Erkennungen profitieren von der integrierten Infrarot-Beleuchtung der Kamera, die Kennzeichen auch bei vollständiger Dunkelheit ausreichend beleuchtet und lesbar macht.
\end{itemize}

\newpage
\subsubsection{Implementierte Maßnahmen zur Datenoptimierung}

Zur Sicherung der Datenqualität im Praxisbetrieb wurden zwei zentrale Strategien angewandt, um Fehlinterpretationen und Datenredundanz zu minimieren:

\paragraph{Zeitfensterbasierte Duplikatserkennung}
Wie in Abschnitt~\ref{sec:duplikatpruefung} beschrieben, verhindert ein konfigurierbarer Intervall-Filter Mehrfacherkennungen desselben Fahrzeugs (z. B. bei langsamer Durchfahrt oder kurzem Stopp im Erkennungsbereich). Dies stellt sicher, dass ein Fahrzeug innerhalb eines logischen Zeitrahmens nur einmal gewertet wird.

\paragraph{Regionale Einschränkung (Ländercodes)}
Die ALPR-Analyse wurde gezielt auf die vier relevanten Länder (AT, HU, SI, DE) beschränkt. Dies verbessert die Genauigkeit signifikant, da das Modell länderspezifische Kennzeichenformate priorisiert und so die Wahrscheinlichkeit von Verwechslungen bei ähnlichen Schriftarten verringert.

Durch diese kombinierten Maßnahmen wird eine für den statistischen Anwendungszweck ausreichend hohe Datenqualität erreicht, welche die beschriebenen physikalischen Limitationen der optischen Erkennung effektiv kompensiert.